Los resultados demuestran que la notación asintótica teórica es insuficiente para dictaminar el desempeño en plataformas de cómputo reales. En ordenamiento, la topología inicial del arreglo condiciona la simetría de partición en \texttt{QuickSort} y el número de pilas en \texttt{PatienceSort}, demostrando la robustez de algoritmos con cotas asintóticas garantizadas como \texttt{MergeSort} e híbridos introspectivos como \texttt{std::sort}. En multiplicación matricial, la ventaja teórica de Strassen ($\mathcal{O}(N^{2.81})$) es neutralizada para dimensiones prácticas ($N \le 1024$) debido al costo de gestión de memoria dinámica en el \textit{heap} y la degradación de la localidad de caché frente a un triple bucle iterativo contiguo vectorizado por el compiladorhola.